\subsection{Radial Symmetry}\label{sec:radial-symmetry}

In a 1979 paper, Gidas, Ni and Nirenberg~\cite{GNN} establish radial symmetry of positive solutions to certain nonlinear elliptic equations. The technique is based on the maximum principle. In this section, we shall consider the simplest case of their result, which is the following theorem.

\begin{Theorem}
\label{thm:radial-symmetry-positive-solution}
Let $f:\mathbb{R}\to\mathbb{R}$ be of class $C^1$. In the ball $B=B_R(0)\subset\mathbb{R}^n$, let $u>0$ be a solution (of class $C^2(\overline B)$) of the equation
\begin{equation}\label{eq:radial-symmetry-semilinear-equation}
\Delta u+f(u)=0,
\end{equation}
satisfying the boundary condition
\begin{equation}\label{eq:radial-symmetry-zero-boundary}
u=0\quad\text{on }\partial B.
\end{equation}
Then $u$ is radially symmetric and strictly monotone decreasing, i.e.,
\begin{equation}\label{eq:radial-symmetry-strict-decrease}
\dfrac{\partial u}{\partial r}<0\quad\text{for }0<r<R.
\end{equation}

\end{Theorem}

The remarkable fact is that the result holds regardless of what $f$ is.

\subsubsection{Two Auxiliary Lemmas}\label{subsec:radial-auxiliary-lemmas}

Throughout the rest of this section, let $B$ and $u$ be as given by the assumptions of Theorem~\ref{thm:radial-symmetry-positive-solution}.

\begin{Lemma}
\label{lem:radial-boundary-directional-derivative}
Let $\bm{x}_0\in\partial B\cap\{x_1>0\}$. Then, for some $\delta>0$, we have $\partial u/\partial x_1<0$ in $B\cap\{|\bm{x}-\bm{x}_0|<\delta\}$.
\end{Lemma}

\begin{Proof}
To simplify notation, we shall write $u_1$ for $\partial u/\partial x_1$, $u_{11}$ for $\partial^2u/\partial x_1^2$, etc.

If the lemma were false, there would be a sequence of points $\bm{x}_j\in B$ with $u_1(\bm{x}_j)\geq0$ and $\bm{x}_j\to\bm{x}_0$. On the other hand, since $u>0$ in $B$ and $u=0$ on the boundary, we must have $u_1(\bm{x}_0)\leq0$. Hence, by continuity, we find $u_1(\bm{x}_0)=0$. Note that since $u=0$ on $\partial B$ the tangential derivatives are all zero. Since the $x_1$ direction is out of the tangent plane, we have shown that $\nabla u(\bm{x}_0)=0$.

We claim that the second derivative $u_{11}(\bm{x}_0)$ is also zero. First, since $u(\bm{x}_0)=u_1(\bm{x}_0)=0$, and $u>0$ in $B$, we must have $u_{11}(\bm{x}_0)\geq0$. Assume now that $u_{11}(\bm{x}_0)>0$. Then $u_{11}$ is also positive in a neighborhood $N$ of $\bm{x}_0$. Consider now the straight line segment $\Gamma$ in the positive $x_1$ direction connecting $\bm{x}_j$ to a point $\bm{y}_j$ on $\partial B$. For sufficiently large $j$, $\Gamma$ lies completely in $N$ so that
\begin{equation}\label{eq:radial-boundary-second-derivative-segment}
u_1(\bm{y}_j)-u_1(\bm{x}_j)=\int_\Gamma u_{11}(\bm{x})\,\mathrm{d}x_1
\end{equation}
is positive. But this is a contradiction since $u_1(\bm{x}_j)\geq0$ and $u_1(\bm{y}_j)\leq0$.

Suppose now that $f(0)\geq0$. Then we have
\begin{equation}\label{eq:radial-nonlinearity-zero-value-shift}
\Delta u+f(u)-f(0)\leq0,
\end{equation}
and by the mean value theorem, we can find a function $c(\bm{x})$ so that $f(u)-f(0)=c(\bm{x})u$. Now Lemma~\ref{lem:elliptic-boundary-point}, applied to $-u$, implies $u_1(\bm{x}_0)<0$, a contradiction.

Hence suppose $f(0)<0$. Then $\Delta u(\bm{x}_0)=-f(0)>0$. On the other hand, since $u=0$ on $\partial B$ and $\nabla u(\bm{x}_0)=0$, it is easy to check that $u_{11}(\bm{x}_0)=n_1^2\Delta u(\bm{x}_0)\ne0$, where $\bm{n}$ denotes the unit normal to $\partial B$ (see Problem~\ref{ex:boundary-hessian-radial-identity}). Again we have a contradiction.
\end{Proof}

For $\lambda\in\mathbb{R}$, let $T_\lambda$ denote the plane $x_1=\lambda$ and let $\Sigma(\lambda)=B\cap\{x_1>\lambda\}$. Moreover, let $\Sigma^{\prime}(\lambda)$ denote the reflection of $\Sigma(\lambda)$ across $T_\lambda$, and let $\bm{x}^\lambda$ denote the reflection of a point $\bm{x}$ across $T_\lambda$. We have the following lemma.

\begin{Lemma}
\label{lem:radial-reflection-strictness}
Assume that for some $\lambda\in[0,R)$ we have
\begin{equation}\label{eq:radial-reflection-weak-comparison}
u_1(\bm{x})<0,\quad u(\bm{x})\leq u(\bm{x}^\lambda)\quad\forall\bm{x}\in\Sigma(\lambda),
\end{equation}
but that $u(\bm{x})$ does not identically equal $u(\bm{x}^\lambda)$ in $\Sigma(\lambda)$. Then $u(\bm{x})<u(\bm{x}^\lambda)$ in all of $\Sigma(\lambda)$ and $u_1(\bm{x})<0$ on $B\cap T_\lambda$.
\end{Lemma}

\begin{Proof}
In $\Sigma^{\prime}(\lambda)$, let $v(\bm{x})=u(\bm{x}^\lambda)$. Then $v$ satisfies the equation $\Delta v+f(v)=0$. Let $w=v-u$, and let $c(\bm{x})$ be such that $f(v)-f(u)=c(\bm{x})w$ (mean value theorem). Then we have
\begin{equation}\label{eq:radial-reflection-difference-equation}
\Delta w+c(\bm{x})w=0
\end{equation}
in $\Sigma^{\prime}(\lambda)$. Moreover, by the assumptions of the lemma we have $w\leq0$ in $\Sigma^{\prime}(\lambda)$ and $w$ is not identically zero in $\Sigma^{\prime}(\lambda)$. Moreover, $w$ vanishes on $T_\lambda\cap B$, which is part of the boundary of $\Sigma^{\prime}(\lambda)$. It now follows from Theorem~\ref{thm:strong-maximum-principle} and Lemma~\ref{lem:elliptic-boundary-point} that $w<0$ in $\Sigma^{\prime}(\lambda)$ and $w_1>0$ on $T_\lambda\cap B$. Since $w_1=-2u_1$ on $T_\lambda\cap B$, the lemma follows.
\end{Proof}

\subsubsection{Proof of the Theorem}\label{subsec:radial-symmetry-theorem-proof}

We shall use the two lemmas from the previous subsection to show that $u$ is symmetric with respect to the plane $x_1=0$ and that $u_1<0$ for $\bm{x}\in B\cap\{x_1>0\}$. This obviously implies the theorem, because by the same argument, $u$ is symmetric with respect to any plane through the origin, and hence radially symmetric. We shall show the following.

\begin{Lemma}
\label{lem:moving-plane-comparison}
For any $\lambda\in(0,R)$, we have
\begin{equation}\label{eq:moving-plane-strict-comparison}
u_1(\bm{x})<0,\quad u(\bm{x})<u(\bm{x}^\lambda)\quad\forall\bm{x}\in\Sigma(\lambda).
\end{equation}

\end{Lemma}

By continuity, we obtain $u(\bm{x})\leq u(\bm{x}^0)$ in $\Sigma(0)=B\cap\{x_1>0\}$. Repeating the same argument with $x_1$ replaced by $-x_1$, we find that $u(\bm{x})\geq u(\bm{x}^0)$ in $\Sigma(0)$. Hence $u$ is symmetric across the plane $x_1=0$, and the theorem follows.

\begin{Proof}
From Lemma~\ref{lem:radial-boundary-directional-derivative}, it follows that \eqref{eq:moving-plane-strict-comparison} holds for $\lambda$ sufficiently close to $R$. Now let $\mu$ be a critical value such that \eqref{eq:moving-plane-strict-comparison} holds for $\lambda>\mu$, but not beyond. By continuity, we find that
\begin{equation}\label{eq:moving-plane-critical-comparison}
u_1(\bm{x})<0,\quad u(\bm{x})\leq u(\bm{x}^\mu)\quad\forall\bm{x}\in\Sigma(\mu).
\end{equation}
We need to show that $\mu=0$. Assume the contrary, i.e., $\mu>0$. For any point $\bm{x}_0\in\partial\Sigma(\mu)\setminus T_\mu$, we have $\bm{x}_0^\mu\in B$, and hence $0=u(\bm{x}_0)\ne u(\bm{x}_0^\mu)$. Hence $u(\bm{x})$ does not identically equal $u(\bm{x}^\mu)$ in $\Sigma(\mu)$, and Lemma~\ref{lem:radial-reflection-strictness} is applicable. Thus $u(\bm{x})<u(\bm{x}^\mu)$ in $\Sigma(\mu)$ and $u_1<0$ on $B\cap T_\mu$. Hence \eqref{eq:moving-plane-strict-comparison} holds for $\lambda=\mu$. Moreover, by Lemma~\ref{lem:radial-boundary-directional-derivative}, we have $u_1<0$ in a neighborhood of any point on $T_\mu\cap\partial B$. Thus every point on $T_\mu\cap\overline B$ has a neighborhood where $u_1<0$, and since $T_\mu\cap\overline B$ is compact, we must have $u_1<0$ in $\Sigma(\mu-\epsilon)$ for $\epsilon$ sufficiently small.

Since we assumed that \eqref{eq:moving-plane-strict-comparison} does not hold beyond $\mu$, there must be a sequence $\lambda_j$ and $\bm{x}_j\in\Sigma(\lambda_j)$ such that $\lambda_j\to\mu$ and
\begin{equation}\label{eq:moving-plane-counterexample-sequence}
u(\bm{x}_j)\geq u(\bm{x}_j^{\lambda_j}).
\end{equation}
After taking a subsequence, we may assume that $\bm{x}_j$ converges; the limit $\bm{x}$ is necessarily in the closure of $\Sigma(\mu)$. Passing to the limit in \eqref{eq:moving-plane-counterexample-sequence}, we find $u(\bm{x})\geq u(\bm{x}^\mu)$. This cannot hold if $\bm{x}\in\partial B$, since in that case $u(\bm{x})=0$, $u(\bm{x}^\mu)>0$. Since \eqref{eq:moving-plane-strict-comparison} holds for $\mu$, we must have $\bm{x}\in T_\mu\cap B$; hence $\bm{x}^\mu=\bm{x}$. On the other hand, the straight line segment connecting $\bm{x}_j$ to $\bm{x}_j^{\lambda_j}$ belongs to $B$, and because of \eqref{eq:moving-plane-counterexample-sequence} and the mean value theorem, it must contain a point $\bm{y}_j$ where $u_1(\bm{y}_j)\geq0$. Taking the limit $j\to\infty$, we find $u_1(\bm{x})\geq0$. This is a contradiction.
\end{Proof}

\paragraph*{Problems}

\begin{enumerate}[label=\thesection.\arabic*.,ref=\thesection.\arabic*,start=17,series=book-problems]

\item\label{ex:nonnegative-nonlinearity-positive-solution} Show that, if $f(u)\geq0$ for every $u$, then nontrivial solutions of \eqref{eq:radial-symmetry-semilinear-equation}, \eqref{eq:radial-symmetry-zero-boundary} are automatically positive. Also show that if $f(0)\geq0$, then \eqref{eq:radial-symmetry-strict-decrease} holds also at $r=R$.

\item\label{ex:radial-center-nonlinearity-positive} Show that if a non-negative solution $u$ of \eqref{eq:radial-symmetry-semilinear-equation}, \eqref{eq:radial-symmetry-zero-boundary} exists and $u$ is not identically zero, then $f(u(0))$ must be strictly positive.

\item\label{ex:boundary-hessian-radial-identity} Let $B$ be a ball centered at the origin. Let $u$ be a $C^2$-function such that $u>0$ in $B$ and $u=0$ on $\partial B$. Let $\bm{x}$ be a point on $\partial B$.

(a) Show that $\Delta u(\bm{x})=u_{rr}(\bm{x})+\frac{n-1}{r}u_r(\bm{x})$ where $r$ is the radial direction.

(b) Assume that $u_r(\bm{x})=0$. Show that $\partial^2u/\partial x_i\partial x_j(\bm{x})=n_in_ju_{rr}(\bm{x})$. Here $\bm{n}$ is the unit normal to $\partial B$. Hint: Show first that the matrix of second derivatives is positive semidefinite.

\item\label{ex:radial-profile-semilinear-construction} Let $u\in C^4(\overline B)$ be radially symmetric and satisfy $u_r<0$ for $0<r\leq R$, and assume that $\Delta u(0)\ne0$. Show that there is a $C^1$-function $f$ such that $\Delta u+f(u)=0$.

\item\label{ex:annular-positive-solution-outward-monotonicity} Let $u>0$ be a $C^2$-solution of \eqref{eq:radial-symmetry-semilinear-equation} in the domain $R_1<|\bm{x}|<R_2$. Assume that $u=0$ on the outer boundary $|\bm{x}|=R_2$. Show that $u_r<0$ for $(R_1+R_2)/2\leq|\bm{x}|<R_2$.

\end{enumerate}
